\documentclass[../main]{subfiles}
\begin{document}
\ifSubfilesClassLoaded{\setcounter{page}{7}}{}
\section{Preface}
\label{sec:02_preface}

In his Critique of the Gotha Programme, Marx wrote: “What we have to deal with here is a communist society, not as it has developed on its own foundations, but, on the contrary, just as it emerges from capitalist society; which is thus in every respect, economically, morally and intellectually, still stamped with the birthmarks of the old society from whose womb it emerges”, and “In a higher phase of communist society, after the enslaving subordination of the individual to the division of labour, and therewith also the antithesis between mental and physical labour, has vanished; after labour has become not only a means of life but life’s prime want; after the productive forces have also increased with the all-round development of the individual, and all the springs of social wealth flow more abundantly—only then can the narrow horizon of bourgeois right be crossed in its entirety and society inscribe on its banners: From each according to his ability, to each according to his needs!”

The idea of private property and the rule of capital is opposed by the idea of communism. In reality, the only thing that can be set against a society in which capital rules is a communist society. But what is communism?

In this pamphlet we are concerned with what Marx called the higher phase of communism, that is, a communist society that has developed on its own foundations, or communism as such. The lower phase of communism, or the period of transition from capitalism to communism, which later came historically to be known as “socialism”, will be examined later. For now, the question is what exactly this transition is supposed to lead to, or what communists want to obtain as a result of the revolutionary transformation of social life.

What interests us here are the essential determinations that make a society specifically communist and distinguish it from all the modes of production and ways of conducting social life that preceded communism. We will therefore set aside any description of technology under communism. It may take various forms (and will, of course, change rapidly). For example, capitalism in the age of the steam engine and capitalism with modern digital technology differ enormously from one another. But what makes capitalism capitalism—a society of developed commodity production in which labour power, among other things, is a commodity, and capital is the dominant social relation—is characteristic, necessary and invariable for capitalism in general, in the 19th century and the 21st alike. An understanding of these necessary moments in their necessary interconnection is what constitutes the concept of capitalism.

In exactly the same way, the concept of communism, like that of any other object, includes only what is necessary and sufficient, in its necessary transitions into one another, for a society to be communist.

The difficulty lies in the fact that we are dealing with an object that has not yet become a reality. History has not yet known (at least on our planet) a communist society that has developed on its own foundations (primitive communism is irrelevant in this case).

However, communism long ago became a science. The entire course of the historical development of thought gave rise to the idea of communism, while the course of social life, at the very least, constantly demonstrates the limit of capitalism and the need to move to a new foundation of social development. Capitalist crises of overproduction are accompanied by the colossal destruction not only of the commodities produced or of individual enterprises, but often of entire once-thriving megacities.

If we consider social production as a mode of interaction between humans and nature, it becomes obvious that it is irrational and is inflicting harm on our planet on such a scale that the damage may become irreparable. The production and overproduction of that same packaging, which is excessive and pointless from the standpoint of a rational interaction with nature but necessary for the commodity to be sold, does enormous harm. The production of commodities themselves, which neither enrich nor develop people, and often even cripple them, yet consume vast quantities of resources and pollute the soil, water and air—this is the reality of capitalism. There is no well-thought-out waste recycling planned on the scale of society as a whole, nor any rational distribution and use of all natural resources and productive capacities. All of this creates the threat of an irreversible ecological catastrophe.

Alongside this man-made destruction of the natural world around us comes the degradation of humanity itself. This is not only the absolute poverty into which a substantial portion of the world's working people have been plunged: people go hungry, live in unsanitary conditions, and have no access to education or healthcare. It is also a matter of relative impoverishment, whereby new needs are produced whose satisfaction becomes a necessity even just to reproduce oneself as a seller of one's own labour power, yet the means to satisfy them are insufficient. Engels already noted this tendency, but today, when capitalism cannot survive without purposefully producing needs, it is intensifying many times over and assuming alarming proportions, especially in the periphery and semi-periphery of the capitalist world, where, on a societal scale, labour power is not fully restored. This undermines the very foundations of capitalist production.

But that is not all. In order to reproduce itself on an expanded scale, modern capital revives and subordinates to itself pre-capitalist forms of the exploitation of man by man. Slavery and serfdom exist in the modern world insofar as they prove profitable and contribute to the self-expansion of value. For example, according to ILO data published in 2014, profits from the use of slave labour worldwide exceed 150.2 billion dollars (as you can see, the profit from slavery is calculated, and can be calculated, in purely capitalist terms rather than in kind, as it was in the era of slaveholding).

Capital constantly demands room for self-expansion, and to achieve this it must push other capital out of existing markets and redistribute resources in its own favour. This generates conflicts not only within states but, above all, between them, which gives rise to local military actions and leads towards yet another world war. And when we have a tendency towards world war (its preparation, its unfolding, its individual flare-ups) combined with nuclear weapons, what accidents are not possible!… At the same time, war itself gives a powerful impetus to the development of a broad, mass, spontaneous anti-capitalist movement. How consistent this line will be, and whether it will lead to a genuine way beyond capitalism, is hard to say for now. In general, no one today can guarantee that humanity will outlive capitalism. The only thing this testifies to is the limit of development of a society in which capital dominates.

That capitalism will come to an end is beyond doubt. The only question is whether it will end together with humanity, or whether humanity will learn to live differently. And how exactly can we live differently, so that humanity as a whole and all its diverse components develop not by hindering but by helping one another? The scientific theory of communism provides the answer to this question, which is by no means outdated but, on the contrary, burningly relevant for humanity.

It is the main points of this theory that we will try to set out here. It was developed by Marx and Engels. Later Marxist thinkers also made their contribution to the development of the understanding of communism. Among them, Lenin's theoretical work deserves particular attention.

The meanings of the terms “communism” and “socialism” did not become established at once. For a long time they were used as synonyms, and communism was sometimes even understood as preceding socialism. Therefore, one should turn equally both to those works of Marx and Engels in which the term “communism” is used and to those that speak of “socialism”, provided it is clear from the substance of the text that what is being discussed is communist society, the communist movement, or the theory of communism. As for Lenin and other Marxists of the 20th century, in their works, as in Marx’s later works, the terminology had already become more or less settled. In his article “A Great Beginning”, Lenin specifically reminds readers that “the scientific distinction between socialism and communism is only that the first term means the first stage of the new society arising out of capitalism, while the second means the higher, further stage of it”. Wherever we use these words without references or special explanations, they should be understood according to contemporary usage.

As for the argument against communism that the various ideological servants of today’s capitalist ruling class keep repeating in every possible way—“communism is a utopia; people have already tried to build communism, and nothing (or nothing good) came of it”—our answer is as follows. Yes, it has not yet been possible to build a communist society, but to say that nothing good came of the attempt is simply a lie. In their day, the successes of the socialist Soviet Union could not be denied even by the most zealous anti-communists, even in the USA. Those who remember the recent history of the 20th century know very well that society’s movement towards communism enabled all the peoples living in the USSR to rise to heights never before seen under feudalism or capitalism, and to do so in the shortest possible time. This is borne out both by dry figures and calculations and by achievements in culture that remain unsurpassed to this day.

If we compare the first decades\footnote{This was despite the fact that the movement began in a war-ravaged, largely agrarian country surrounded by hostile powers; it was a movement along an untrodden path and therefore proceeded amid great problems, difficulties and retreats, and even tragedies, which we communists have no right to forget. Moreover, this movement took place in a society with rather strong capitalist tendencies, which cannot be abolished by decree or simply stepped over without doing away with commodity production.} in the life of a society moving towards communism with the same span of time in which the “independent states” have been going their separate ways since the collapse of the USSR, we find a direct correlation between movement towards communism and the development of society as a whole. This concerns not only, or even primarily, industry and the production of goods, but rather the opportunities for the development of the individual and of human abilities available to the broad masses of the population.

As early as 1925, there were 2,250,000 more pupils in primary schools than in 1913. The number of students in vocational schools doubled. Society went straight from the Civil War to the classroom. The right to an education was won for those who, before the Civil War, had been doomed to ignorance. Just think about it: the Soviet Union was only a year old, and it already had something to show the capitalist world, where the estate principle dominated education. In the capitalist world, a good education was, and still is, out of reach for the majority.

Within 20 years of October, the problem of child homelessness was solved, and the children who had lost their parents during the First World War and the Civil War were given not just a roof over their heads and bread, but an education, a trade, a start in life. Meanwhile, homeless children have reappeared amid the fragments of what was once a single socialist country—and most of them have living parents. The general level of education of the population has fallen and continues to fall. The education systems of the former Soviet republics are striving to conform to the new social realities: ever-narrower specialisation and a deepening divide between the educational opportunities available to the children of the poor and those of the rich. You will find a wealth of telling evidence of this process in Russia in the film “The Last Bell” by Konstantin Syomin and Yevgeny Spitsyn.

In the USSR, the campaign to eliminate illiteracy (likbez, launched in 1920) and the programme of universal primary education (vseobuch, from 1930) managed to solve the problem of primary education in less than twenty years—and this in an agrarian country where, before 1917, more than 80\% of the population was illiterate or semi-literate. By 1936, around 40 million illiterate people had been taught to read and write in the USSR. Between 1933 and 1937, more than 20 million illiterate and about 20 million semi-literate people studied in the officially registered likbez schools alone. According to the 1939 census, literacy among people aged 16 to 50 was approaching 90\%. It was these very people—not mythical warrior-heroes with crosses on their chests and God in their hearts (as Nikita Mikhalkov’s film “The Citadel” and other products of today’s propaganda industry would have it)—who defeated the war machine of Nazi Germany, the most powerful, modern and well-equipped of its day, which by the time of the invasion of the USSR had most of European industry working for it.

So, within 20 years, illiteracy was virtually eliminated and a network of medical facilities was built that gave everyone access to healthcare (doctors and teachers appeared in places that had never had them before); schools, libraries and cultural institutions sprang up, and science received a powerful boost, thanks in part to the creation of the necessary infrastructure. The building of communism, understood as prioritising the all-round development of the individual (education, healthcare, science, culture), made it possible to defeat fascism, which threatened the entire planet, and then, after a hard and devastating war, to rebuild the country in record time.

And what did life in the years after the breakup of the USSR—the complete dismantling of communist tendencies in social development—bring to the former Soviet republics?

In Ukraine alone, according to the State Statistics Committee, the absolute population decline from the collapse of the USSR up to 2014 amounted to almost 6.5 million people (before the collapse, population growth had been positive). This figure does not account for those who are still “registered” as citizens of Ukraine but have been forced to leave their homes to find work elsewhere. This in itself testifies to a colossal social catastrophe. Of course, someone will bring up the famine of 1932–1933. Mass famine is always a terrible tragedy, regardless of its true causes. But even according to the most anti-Soviet calculations (calculations, mind you, not groundless ideological claims) carried out by the Institute for Demography of the National Academy of Sciences of Ukraine, the famine in Ukraine claimed not 7 or 10 million lives, but 3.9 million. These calculations did not simply count losses; they compared the actual population with the projected population, taking into account the significant positive growth in the years before and after the famine. If the absolute population losses in Ukraine after the breakup of the USSR were calculated using the same method, the picture would be far more horrifying than the nearly 6.5 million obtained by simple subtraction.

Lithuania, Latvia and Estonia have also recorded significant depopulation. In the Russian Federation, thanks to migration, the decline was smaller: from 148.6 million in 1991 to 144.3 million in 2016.

If we take into account labour migration in other directions and the employment structure within the Central Asian republics themselves, it becomes clear that the overall quality of life, and the chance for citizens of these republics to find work at home that can support a family, have declined considerably.

Wars on the territory of the former USSR have become a frequent and commonplace occurrence: the conflict between Azerbaijan and Armenia in Nagorno-Karabakh (which began while the USSR still existed, but essentially during its collapse), the Civil War in Tajikistan, Transnistria, the two Chechen wars, Georgia and Abkhazia, South Ossetia, Donbass. And in all these wars people killed one another, in the last instance, to serve the interests of rival factions of capital—that is, their own exploiters.

Some may object that the Civil War that began in 1917 was also bloody and claimed, by various estimates, between 8 and 13 million lives on all sides. This figure includes not only those killed in combat but also those who died of disease, famine and terror. Yes, that is true, we reply. And, of course, it would have been better to avoid war altogether. But, first, it was not the Bolsheviks or the proletariat who unleashed it, but representatives of the old exploiting classes who did not want to lose their dominance. Second, in that war the working people were defending their own interests, not the interests of those who live off the exploitation of their labour. Third, that war was a struggle for the opportunity to carry out a radical transformation of the whole of society, not merely a redivision of capital’s spheres of influence. The Great Patriotic War deserves to be discussed separately, since it was a struggle against fascism in which the Soviet people defended not only their own interests but those of all humanity.

As you can see, even such a simple comparison of the first decades of the USSR with the life of the republics after its collapse allows us to speak of the advantages of a society in which there is room for the tendency towards communism. And this despite the fact that the process was far from smooth, and there was much that could by no means be called communist. We are referring to purely capitalist processes and phenomena within socialist society. It is precisely these that critics of the USSR make so much noise about, even though this criticism, where it is entirely justified, turns out to be self-criticism with respect to the capitalist world.

The tendency towards communism in the USSR made it possible, just twelve years after the Great Patriotic War, which claimed millions of lives and destroyed a substantial part of the country, to launch the first artificial Earth satellite—and soon afterwards to send the first human into space. The building of communism made it possible to work out scientifically ideas that, for us in the present state of society, are ideas of the future rather than of the past, since their realisation is incompatible with private property in the means of production. The list could go on at length, and it could include not only Soviet achievements but also those of other socialist countries—countries that arose precisely along the path Lenin envisioned in the light of the hypothesis of a victorious world revolution.

Of course, this grand construction project involved not only achievements but also retreats and failures. But even if we set aside what has been said above, the argument

“we failed to build communism, so it is impossible altogether and nothing will ever come of it” suffers from the same logical flaw as the claims made by opponents of the ideas that underpinned, for example, the emergence and development of aviation. Those opponents insisted that it was impossible to build flying machines heavier than air. It is the same as saying that a person, or humanity as a whole, will never be able to learn something simply because they have not managed it so far, even though attempts have been made.

Such illogical arguments involve no examination of the essence of the matter and no substantive objections. That is why they have nothing whatsoever to do with the science of communism. The failure of humanity’s first leap towards communism testifies only to the need for a detailed study of historical experience in order to understand the reason for that failure, and for a more thorough study of the conditions that must be achieved for the transition to communism. The very existence and content of the communist ideal, which was born as an understanding of the tendencies and contradictions of capitalist society, demands this.

So, since communism, among other things, has become a science, let us proceed to study it. This pamphlet, of course, cannot claim to be a detailed and comprehensive study of communism. Rather, it is a text intended to set out the basic tenets of the science of communism as developed in Marxist theory. A deeper and more detailed treatment of the problem is a task for the near future. We invite our readers to join in this work, including in the form of criticism of our own.
\end{document}
